% Options for packages loaded elsewhere
\PassOptionsToPackage{unicode}{hyperref}
\PassOptionsToPackage{hyphens}{url}
\PassOptionsToPackage{expansion=false}{microtype}
%
\documentclass[
]{article}
\usepackage{amsmath,amssymb}
\usepackage{iftex}
\ifPDFTeX
  \usepackage[T1]{fontenc}
  \usepackage[utf8]{inputenc}
  \usepackage{textcomp} % provide euro and other symbols
\else % if luatex or xetex
  \usepackage{unicode-math} % this also loads fontspec
  \defaultfontfeatures{Scale=MatchLowercase}
  \defaultfontfeatures[\rmfamily]{Ligatures=TeX,Scale=1}
\fi
% lmodern not available; using default CM fonts
\ifPDFTeX\else
  % xetex/luatex font selection
\fi
% Use upquote if available, for straight quotes in verbatim environments
\IfFileExists{upquote.sty}{\usepackage{upquote}}{}
\IfFileExists{microtype.sty}{% use microtype if available
  \usepackage[]{microtype}
  \UseMicrotypeSet[protrusion]{basicmath} % disable protrusion for tt fonts
}{}
\makeatletter
\@ifundefined{KOMAClassName}{% if non-KOMA class
  \IfFileExists{parskip.sty}{%
    \usepackage{parskip}
  }{% else
    \setlength{\parindent}{0pt}
    \setlength{\parskip}{6pt plus 2pt minus 1pt}}
}{% if KOMA class
  \KOMAoptions{parskip=half}}
\makeatother
\usepackage{xcolor}
\setlength{\emergencystretch}{3em} % prevent overfull lines
\providecommand{\tightlist}{%
  \setlength{\itemsep}{0pt}\setlength{\parskip}{0pt}}
\setcounter{secnumdepth}{-\maxdimen} % remove section numbering
\ifLuaTeX
  \usepackage{selnolig}  % disable illegal ligatures
\fi
\IfFileExists{bookmark.sty}{\usepackage{bookmark}}{\usepackage{hyperref}}
\IfFileExists{xurl.sty}{\usepackage{xurl}}{} % add URL line breaks if available
\urlstyle{same}
\hypersetup{
  hidelinks,
  pdfcreator={LaTeX via pandoc}}

\author{}
\date{}

\begin{document}

\hypertarget{bayes-theorem-a-foundation-of-probabilistic-reasoning}{%
\section{Bayes' Theorem: A Foundation of Probabilistic
Reasoning}\label{bayes-theorem-a-foundation-of-probabilistic-reasoning}}

\hypertarget{abstract}{%
\subsection{Abstract}\label{abstract}}

Bayes' theorem describes how to update beliefs in light of new evidence.
First formulated by Thomas Bayes in the 18th century and later developed
by Pierre-Simon Laplace, it has become a cornerstone of statistics,
machine learning, and scientific inference. This document presents the
theorem, its derivation, key interpretations, and several illustrative
applications including medical testing and spam filtering.

\hypertarget{introduction}{%
\subsection{1. Introduction}\label{introduction}}

Probability theory provides a rigorous language for reasoning under
uncertainty. Among its most powerful tools is Bayes' theorem, which
relates the conditional probability of a hypothesis given observed data
to the prior probability of the hypothesis and the likelihood of the
data. In an era of data-driven decision-making, Bayesian reasoning
underlies applications ranging from clinical trials to autonomous
vehicles.

The central question Bayes' theorem answers is: \emph{given that we have
observed some evidence, how should we rationally revise our belief about
a hypothesis?}

\hypertarget{mathematical-formulation}{%
\subsection{2. Mathematical
Formulation}\label{mathematical-formulation}}

\hypertarget{conditional-probability}{%
\subsubsection{2.1 Conditional
Probability}\label{conditional-probability}}

The conditional probability of event \(A\) given event \(B\) is defined
as:

\[P(A \mid B) = \frac{P(A \cap B)}{P(B)}, \quad P(B) > 0\]

By symmetry, \(P(B \mid A) = \frac{P(A \cap B)}{P(A)}\), so:

\[P(A \cap B) = P(A \mid B)\,P(B) = P(B \mid A)\,P(A)\]

\hypertarget{bayes-theorem}{%
\subsubsection{2.2 Bayes' Theorem}\label{bayes-theorem}}

Rearranging the identity above yields \textbf{Bayes' theorem}:

\[\boxed{P(A \mid B) = \frac{P(B \mid A)\,P(A)}{P(B)}}\]

In the context of hypothesis testing, we write \(H\) for a hypothesis
and \(D\) for observed data:

\[P(H \mid D) = \frac{P(D \mid H)\,P(H)}{P(D)}\]

The terms have well-known names:

\begin{itemize}
\tightlist
\item
  \(P(H)\) --- the \textbf{prior}: our belief in \(H\) before seeing
  \(D\)
\item
  \(P(D \mid H)\) --- the \textbf{likelihood}: how probable the data are
  if \(H\) is true
\item
  \(P(H \mid D)\) --- the \textbf{posterior}: our updated belief after
  seeing \(D\)
\item
  \(P(D)\) --- the \textbf{marginal likelihood} (evidence): a
  normalising constant
\end{itemize}

\hypertarget{the-law-of-total-probability}{%
\subsubsection{2.3 The Law of Total
Probability}\label{the-law-of-total-probability}}

When \(H\) takes discrete values \(H_1, H_2, \ldots, H_n\) that
partition the sample space, the denominator expands via the law of total
probability:

\[P(D) = \sum_{i=1}^{n} P(D \mid H_i)\,P(H_i)\]

This gives the full form:

\[P(H_k \mid D) = \frac{P(D \mid H_k)\,P(H_k)}{\displaystyle\sum_{i=1}^{n} P(D \mid H_i)\,P(H_i)}\]

\hypertarget{applications}{%
\subsection{3. Applications}\label{applications}}

\hypertarget{medical-diagnostic-testing}{%
\subsubsection{3.1 Medical Diagnostic
Testing}\label{medical-diagnostic-testing}}

Consider a disease affecting 1\% of a population. A test has 99\%
sensitivity (true-positive rate) and 95\% specificity (true-negative
rate, so 5\% false-positive rate). Given a positive test result, what is
the probability the patient has the disease?

Let \(H\) = ``patient has disease'' and \(D\) = ``test is positive'':

\[P(H) = 0.01, \quad P(D \mid H) = 0.99, \quad P(D \mid \neg H) = 0.05\]

\[P(D) = 0.99 \times 0.01 + 0.05 \times 0.99 = 0.0099 + 0.0495 = 0.0594\]

\[P(H \mid D) = \frac{0.99 \times 0.01}{0.0594} \approx 0.167\]

Despite the accurate test, only about 16.7\% of positive results
correspond to true disease. This counterintuitive result is a direct
consequence of the low base rate, illustrating why prior probabilities
matter enormously.

\hypertarget{bayesian-spam-filtering}{%
\subsubsection{3.2 Bayesian Spam
Filtering}\label{bayesian-spam-filtering}}

A spam filter estimates \(P(\text{spam} \mid \text{words in email})\).
Using a naive Bayes classifier, it assumes word occurrences are
conditionally independent given the class, then applies:

\[P(\text{spam} \mid w_1, \ldots, w_n) \propto P(\text{spam})\prod_{i=1}^{n}P(w_i \mid \text{spam})\]

Training data provide the likelihoods \(P(w_i \mid \text{spam})\) and
\(P(w_i \mid \text{ham})\), and the prior \(P(\text{spam})\) is
estimated from the proportion of spam in the training set.

\hypertarget{the-bayesian-vs.-frequentist-debate}{%
\subsection{4. The Bayesian vs.~Frequentist
Debate}\label{the-bayesian-vs.-frequentist-debate}}

Frequentists interpret probability as the long-run frequency of events;
they do not assign probabilities to hypotheses. Bayesians treat
probability as a degree of belief, permitting statements like ``there is
a 73\% chance this hypothesis is true.'' The two schools differ in how
they handle prior information:

\begin{itemize}
\tightlist
\item
  \textbf{Frequentists} rely solely on the likelihood function and avoid
  priors.
\item
  \textbf{Bayesians} explicitly encode prior knowledge, which is then
  updated by data.
\end{itemize}

In practice, both approaches are valuable; the Bayesian framework is
especially useful when prior knowledge is available or when sequential
updating of beliefs is required.

\hypertarget{conclusion}{%
\subsection{5. Conclusion}\label{conclusion}}

Bayes' theorem is a deceptively simple equation with profound
implications. It provides a principled mechanism for learning from
evidence: start with a prior belief, observe data, compute the
likelihood, and obtain a posterior belief. Applications span medicine,
finance, natural language processing, robotics, and fundamental science.
As datasets grow and computational power increases, Bayesian methods
continue to gain prominence as a rigorous framework for inference under
uncertainty.

\end{document}
